\chapter{The limits of the method: the Hermite truncation instability}\label{ch:limits}

The method has one limit that governs everything the paper can claim about strong drives. In a non-uniform field the truncated asymmetrically weighted (AW) Hermite system has growing modes that are artefacts of the basis, they grow faster when the resolution is raised, and no integrator, closure or alternative basis tried so far removes them within the cost and conservation constraints of the project. This chapter tells how the instability was found, what the records measure about it, why it can be derived from the structure of the weighted basis, why time integration cannot cure it, which remedies were tested and why each was rejected, and what is left, namely the certified windows. The physical instability that the strong drive also excites, the oscillating two-stream instability, is separated from it at the end.

\section{How the instability was found}\label{sec:found}

The resonant-drive study began at low spatial resolution. With the pump frame and the exact remaps of Sec.~\ref{sec:pump}, runs at $N_x=8$ and $N_n$ of 64, 128 and 256 agreed with each other to the point where the electron distribution lost positivity, and gave the first windows. Without a closure, at $N_x=8$, the Hermite results for $\vq=0.03$ and $0.1$ were converged in $N_n$ to $\wpe t=879$ and $606.5$ (the lane that fixes $t_{\rm res}=304(\vq)^{-0.30}$ for $0.02\le\vq\le0.1$, an empirical law for those seeds and that box). Raising the spatial resolution did not extend these windows; it shortened them. At $\vq=0.1$ without closure $N_x=16$ failed at $t=476$, while $N_x=8$ reached 700 or more, and raising $N_n$ did not help. The first interpretation was physical. The drive amplifies high-$k$ fields strongly, the exact linear Volterra solution confirms that finite-$k$ modes grow faster the larger $k$ is, and it seemed natural that a finer grid simply resolved a real cascade that a coarser grid could not carry. This reading was written into the results and later retracted. A correction dated 2026-10-19 records that the GPU resolution study refutes it: with a field-scaled closure of strength $c=4$ more resolution ends the run sooner, at $t=595$, 383, 453 and 328 for $(N_x,N_n)=(64,128)$, $(64,256)$, $(128,128)$ and $(128,256)$ at $\vq=0.1$, whereas a physical cascade resolved better should not terminate earlier. The runs in that study all ended at the adaptive step budget with the Hermite spectrum piling up in the top third of the ladder.

The diagnosis then proceeded through a series of exclusions that were each declared first. Aliasing was excluded: the padded convolution equals the $2/3$-masked right-hand side to $10^{-16}$. Step size was excluded: a fixed step $\mathrm dt=0.01$ fails at the same time. The remap was excluded: with remaps off the run fails at 472. The integrator was excluded by the analysis of Sec.~\ref{sec:integrators}. What remained was the operator itself.

\section{The frozen-field operator and its measured spectrum}\label{sec:frozen}

For one species and one velocity dimension, freeze the field at a prescribed profile $E(x)$ and keep Fourier modes $|k|\le K k_1$ and Hermite indices $n<N_n$. The semi-discrete system is $\dot C=(L_s+L_E+L_\nu)C$, with the streaming $L_s=-\ii k\,aT$ (plus a scalar shift by the centre), $T$ the real symmetric tridiagonal Hermite matrix with $T_{n,n+1}=\sqrt{(n+1)/2}$, the force $L_E$ of Eq.~\eqref{eq:hier} acting from index $n-1$ to $n$ as a convolution in $k$, and the diagonal closure $L_\nu$. The streaming is skew-Hermitian, so its exponential is unitary and its eigenvalues $\pm\ii k a\xi_j$ are imaginary, with $\xi_j$ the Gauss--Hermite nodes and the largest of order $\sqrt{2N_n}$. The force is strictly lower triangular in $n$ for a frozen field and hence nilpotent: it has no eigenvalues of its own. Growth therefore requires both parts together, in a non-normal product.

The records show the instability in this operator. The frozen-field spectrum for $a=\sqrt{2\times10^{-3}}$, $u=0$, a single cosine field of peak $E_1$, is real and positive at the top of the ladder for the AW basis and zero to round-off for the symmetric-weight (SW) basis, in which the functions have weight $e^{-\xi^2/2}$ instead of $e^{-\xi^2}$ (Table~\ref{tab:aw}). The largest growth rate rises with $N_n$ at fixed $E_1$ (at $E_1=10^{-2}$ and $K=5$: $0.47$, $0.76$, $1.6$ for $N_n=32$, 64, 128, and $K=5$, $N_n=128$ is the number $1.6$ quoted in the paper) and, at fixed $N_n$, roughly in proportion to a power of $E_1$ slightly above one. In a pump-modulated field the Floquet exponent per period, which is the quantity that matters for the real run, is zero to round-off at the lowest field and $0.013$ to $0.15$ at $E_1=2\times10^{-2}$ for $K=2$--5, again with the SW exponent at $10^{-15}$. The eigenvalues of the AW operator are poorly conditioned. The eigenvector matrix has condition number $10^{19}$ to $10^{29}$, against 40 to 105 for the normal (square-root SW) form of Sec.~\ref{sec:remedies}, and the largest real part for $E_1=10^{-2}$, $K=5$, $N_n=64$ was $0.756$ from a dense eigensolver, $0.846$ from ARPACK and $0.702$ from the asymptotic growth of $e^{Lt}$ applied to a random vector. This is why the later closure study measured time-domain growth and not eigenvalues.

\begin{table}[t]
\caption{Largest real part of the frozen-field spectrum for the AW and SW bases (peak field $E_1$, Fourier cut-off $K$, Hermite order $N_n$) and Floquet exponent in a pump-modulated field (\texttt{studies/highk/}). The SW values are round-off at every point tested.}\label{tab:aw}
\small
\resizebox{\ifdim\width>\textwidth\textwidth\else\width\fi}{!}{%
\begin{tabular}{@{}llrcccc@{}}\toprule
basis & $K$ & $N_n$ & $E_1=10^{-3}$ & $E_1=10^{-2}$ & Floquet $8\times10^{-3}$ & Floquet $2\times10^{-2}$\\\midrule
AW & 2 & 32 & $6.8\times10^{-2}$ & $0.26$ & $10^{-15}$ & $1.3\times10^{-2}$\\
AW & 2 & 64 & $0.10$ & $0.73$ & $2.6\times10^{-3}$ & $4.4\times10^{-2}$\\
AW & 5 & 32 & $9.2\times10^{-2}$ & $0.47$ & $3.5\times10^{-6}$ & $4.3\times10^{-2}$\\
AW & 5 & 64 & $0.14$ & $0.76$ & $2.0\times10^{-5}$ & $0.15$\\
AW & 5 & 128 & $0.20$ & $1.6$ & --- & ---\\
AW & 10 & 128 & $0.21$ & $1.7$ & --- & ---\\
SW & 2--10 & 32--128 & $\le8\times10^{-16}$ & $\le6.4\times10^{-15}$ & $\le2.6\times10^{-15}$ & $\le2.2\times10^{-15}$\\\bottomrule
\end{tabular}}
\end{table}

The closure-strength study put numbers on the scaling. Using time-domain growth of $e^{Lt}$ with the closure of Eq.~\eqref{eq:closure}, strengths $c=0,4,16,64,256,1024$, peak fields $10^{-3}$ and $4\times10^{-3}$ (the $k_1$ field at $\wpe t\sim200$--450 in the $\vq=0.1$ runs), $(K,N_n)\in\{2,5,10,21\}\times\{32,64,128\}$ plus $(5,256)$ and $(10,256)$, and single-cosine, flat and $1/j$ spectra, log-linear fits (residual $0.4$--$0.5$ in $\ln\gamma$, a factor $1.6$) give $\gamma(c=0)\sim32\,N_n^{0.57}\hat E^{1.2}$ with no dependence on $K$ (exponent $0.00$), $\gamma(c=4)\sim2.4\,N_n^{0.78}\hat E^{1.1}$, so that more resolution makes $c=4$ worse, and a floor $\gamma(c=1024)\sim11\,K^{0.15}N_n^{0.30}\hat E^{1.26}$, with $\hat E$ the peak field. At peak $4\times10^{-3}$ the $c=4$ growth is $0.04$--$0.46$ and the floor at $c=1024$ is $0.02$--$0.14$, which over $\wpe t=1000$ amplify a $10^{-14}$ seed by $e^{20}$ to $e^{100}$.

\section{Where the growth comes from}\label{sec:mechanism}

The location (top of the ladder), the dependence on $N_n$ and $E$, and the neutrality of the SW basis can all be read from one transformation, which we give here because it turns the numerical finding into a bound. The AW trial functions are $\psi_n=\varphi_ne^{-\xi^2}$ and the symmetric (orthonormal) Hermite functions are $h_n=\pi^{-1/4}\varphi_ne^{-\xi^2/2}$, so $\psi_n=\pi^{1/4}e^{-\xi^2/2}h_n$. Writing $f=e^{-\xi^2/2}g$, the AW coefficients $C_n$ of $f$ are, up to a constant, the symmetric-basis coefficients of $g=fe^{\xi^2/2}$, and since the test functions $\varphi_n$ of Eq.~\eqref{eq:awfun} are the polynomial parts of the same $h_n$, the AW moment equations are exactly the Galerkin projection of the equation for $g$ onto $\mathrm{span}\{h_0,\dots,h_{N-1}\}$. Insert $f=e^{-\xi^2/2}g$ into the Vlasov equation, using $\partial_vf=e^{-\xi^2/2}(\partial_vg-\xi g/a)$:
\begin{equation}
\partial_tg+v\,\partial_xg+\frac qm E\,\partial_vg=\frac{q}{m}\frac{E}{a}\,\xi\,g .
\label{eq:geq}
\end{equation}
The left side is a divergence-free vector field in $(x,v)$ and so a skew operator in $L^2(x,v)$; its Galerkin projection on an orthonormal basis stays skew, whatever the truncation, and that is the algebraic reason why the SW basis, which is Galerkin on $f$ itself, is neutral to round-off (the Hermite matrices are $\xi\to T$ symmetric and $\partial_\xi\to$ antisymmetric tridiagonal). The right side of Eq.~\eqref{eq:geq} is a real multiplication operator, symmetric, which in the retained space is $(q/ma)E(x)\otimes T$ with spectrum bounded by $(|q|/ma)\max_x|E|\,\xi_{\max}$, where the largest Gauss--Hermite node satisfies $\xi_{\max}\simeq\sqrt{2N}$. Since the skew part does not change $\|g\|^2$, the norm of the AW solution in the $g$ variable obeys
\begin{equation}
\frac{\dd}{\dd t}\|g\|\le\frac{|q|}{m}\frac{\sqrt{2N}\,\max_x|E|}{a}\,\|g\|,
\label{eq:growthbound}
\end{equation}
and the multiplication by $\xi$ is largest at the top of the ladder, where the retained eigenvectors of $T$ live at the largest $|\xi|$. This accounts for the location of the unstable modes in the top third of the ladder. The expression in Eq.~\eqref{eq:growthbound} is the natural scale of the field-scaled closure of Eq.~\eqref{eq:closure}, which was obtained by reasoning about the force operator and is now seen to be the exponential growth bound of the non-skew part. Against the records the bound holds everywhere it can be tested and is not tight: at $E_1=10^{-2}$ it gives $1.8$ and $3.6$ for $N_n=32$ and 128 against the measured $0.26$ and $1.6$; at $10^{-3}$ it gives $0.18$ and $0.36$ against $0.068$ and $0.15$--$0.21$; and the measured exponents ($N_n^{0.57}$, $\hat E^{1.2}$) are close to those of the bound ($N_n^{1/2}$, $\hat E^1$). That the bound is not saturated is what one expects from a non-normal operator whose skew part rotates the amplitude out of the maximally growing direction, and the measured rates are $0.15$ to $0.45$ of the bound. We offer this as an explanation consistent with the records, not as a theorem about the observed rates; the design note of the repository reaches the same conclusion about non-normality from the numerical evidence.

The same structure explains why a weighted basis is attractive despite the instability: a Maxwellian is a single mode, moments of order $\le n$ are exact, and conservation of mass, momentum and energy is exact by the argument of Sec.~\ref{sec:hermite}; the price is the non-skew term, which is a pure consequence of putting the Gaussian weight on the trial functions only.

\section{Why integrators cannot fix it}\label{sec:integrators}

A consistent time integrator reproduces the semi-discrete growth as the step goes to zero. If the instability is in $L_s+L_E$ no exact treatment of either block changes it, because the sum is non-normal and its growth rate is a property of the sum. This was tested in a frozen-field toy ($K=5$, $N_n=64$, $E_1=4\times10^{-3}$, $a=\sqrt{2\times10^{-3}}$, $q/m=-1$; streaming spectral radius $2.35$, force norm $0.97$, force block nilpotent to round-off). The exact growth rate is $0.448$. The measured per-step growth of the one-step map is in Table~\ref{tab:toy}: a Lawson-RK4 step that exponentiates the streaming and a Strang step with exact sub-flows return $0.448$ at every step tried up to $h=1$, and $0.42$ and $0.38$ at $h=3$, while RK4 itself becomes unstable at $h\simeq1.2$ because $h\rho$ passes $2.8$. The only schemes that lower the growth are implicit damping schemes, and they lower it by an amount fixed by the step and not by the physics: backward Euler gives $0.43$, $0.34$, $0.085$, $0.045$, $0.015$ at $h=0.01$, $0.05$, $0.2$, $1$, $3$. An L-stable implicit step on the top modes is therefore a hidden, step-dependent closure that stops working as $h\to0$, and would have to be declared as a regulariser. The implicit midpoint rule, which is symmetric and conserves quadratic invariants of skew systems, does not damp: $0.448$, $0.446$, $0.43$, $0.22$, $0.075$ at the same steps, and only steps over the fast frequencies. The design note therefore separates two failures: stiffness (a step restriction from the imaginary spectral radius of $L_s$ or from a strong closure, which exponential integrators remove) and instability (a positive real part of the semi-discrete operator, which they do not).

\begin{table}[t]
\caption{Growth rate per unit time of the one-step map for the frozen toy field (exact $0.448$), from \texttt{studies/imex/}. This check was not gated.}\label{tab:toy}
\small
\resizebox{\ifdim\width>\textwidth\textwidth\else\width\fi}{!}{%
\begin{tabular}{@{}lccccc@{}}\toprule
$h$ & RK4 & Lawson-RK4 & midpoint & backward Euler & Strang, exact sub-flows\\\midrule
0.01 & 0.448 & 0.448 & 0.448 & 0.428 & 0.448\\
0.2 & 0.447 & 0.448 & 0.43 & 0.085 & 0.447\\
1 & 0.388 & 0.447 & 0.217 & 0.045 & 0.441\\
3 & 1.4 & 0.421 & 0.075 & 0.015 & 0.384\\\bottomrule
\end{tabular}}
\end{table}

The question of which of the two limits the adaptive runs was answered by a diagnostic study with gates, hypotheses and four dated revisions committed before the measurements they govern. On the $\vq=0.1$ case ($N_x=64$, $N_n=128$, pump frame, GPU) the exact Jacobian of the right-hand side (11\,272 real unknowns) was linearized at checkpoints, with the Dopri8 stability bound $h_{\rm stab}$ computed from the stability function, and a counterfactual with the closure removed from the right-hand side. For $c=4$, which ends at 578 (595 in the earlier record, since the blow-up is segment-sensitive), the spectral radius of the Jacobian equals the closure rate $\nu_e=c\sqrt{2N_n}\max|E|/a$ to $0.2$--$7\,\%$ for $t\ge500$ and the ratio of the accepted step to $h_{\rm stab}$ rises from $0.08$ at $t=300$ to $0.49$ at $t=576$, with the step collapsing by a factor of $92$; the tolerance ratio of steps at $10^{-10}$ and $10^{-4}$ is $2.5$ to $5.3$, so the run is partly accuracy limited. For $c=64$ ($t_f=760$) the spectral radius equals $\nu_e$ to $0.0\,\%$, real and negative (835.6 and 2318.4 at two checkpoints, equal to the radius of the linear block), the ratio of step to $h_{\rm stab}$ is $0.7$ to $1.9$, and the tolerance ratio is $1.02$ to $1.12$ at five of the six near-failure checkpoints, so the run is stiffness limited; removing the closure from the right-hand side cuts the steps per unit time by $6$--$9$ (548 to 88 at $t=760$). The three declared hypotheses, stiffness limited at every near-failure checkpoint (HA), accuracy limited (HC) and growth realized as predicted (HB), all failed as declared, and the declared verdict was MIXED for both strengths. The reading is: the stiff eigenvalue is the closure itself, and the instability, with maximal real part $2$--$3$, does not control the step. The Lawson factor of Sec.~\ref{sec:lawson} therefore removes the stiffness, as it should (5.0 times fewer steps for $c=64$) and the window does not move (760 by positivity for $c=64$, 579.5 against 578 for $c=4$), which the design note had predicted.

\section{The remedies that were tried}\label{sec:remedies}

Eleven distinct remedies were tried over the life of the project. Each had gates declared before the run, and Table~\ref{tab:remedies} lists them in one place; the text below adds the evidence.

\begin{table}[t]
\caption{Remedies for the truncation instability. All gates were declared before the runs. Window means the certified time of Sec.~\ref{sec:windows}; $\vq=0.1$ unless stated.}\label{tab:remedies}
\small
\resizebox{\ifdim\width>\textwidth\textwidth\else\width\fi}{!}{%
\begin{tabular}{@{}L{3.1cm}L{5.4cm}L{5.8cm}@{}}\toprule
Remedy & Result & Why rejected or kept\\\midrule
Moving basis (pump frame, remaps) & fixed basis 365/210 (at 0.03/0.1) $\to$ 870/612 with the pump gate & removes the quiver, not the instability; remaps off fails at 472\\
Hypercollision, $\nu=1$--2 (lane B) & $\ge1000$ at 0.03; at 0.1 every $\nu\in[0.25,64]$ fails before 1000 & postpones positivity loss; removes 82--92\,\% of the echo at $N=128$\\
Exponential filter and width floor (\#61, lane A) & not run to a comparison & kept open, no records in production\\
Field-scaled closure (\#66) & 0.03 to 942--952, 0.1 to 672 at $N_x=32$ & kept in production as a declared regulariser\\
Resolution-scaled closure strength & no $c$ reaches $\gamma\le10^{-3}$ at 39 of 39 points; stage F ceiling 713--757 & gate Q1 failed; premise false for this family\\
Higher resolution, GPU & 0.1: 595, 383, 453, 328; 0.05: 752, 555, 671, 440 & more resolution ends sooner; gates G1, G4 failed\\
Dai stable AW Galerkin & unstable in AW coordinates ($\mathrm{Re}\,\lambda$ up to 1.1), passes in orthonormal form & AW form fails test 1; orthonormal form is option B\\
Symmetric weight (SW) & neutral but every run blows up at $t=637$--738 & fails tests 3, 4; moments not exact\\
Anti-symmetric square-root SW & neutral; departs from the grid at 391 and 498 & fails G3, G4, earlier than SW\\
Lawson closure & 5.0$\times$ fewer steps at $c=64$; window unchanged & kept (stiffness only)\\
Lawson streaming & 9652 vs 9606 steps, $2.7\times$ cost per step & rejected; step collapse at $c=4$\\\bottomrule
\end{tabular}}
\end{table}

Consider the moving basis first. The pump frame carries the coherent quiver in the centre $u_s(t)$ rather than in the high Hermite indices, and the exact remap changes the width and the residual drift between segments. At $N_x=8$ it moves the windows at $\vq=0.03$ and $0.1$ from $365$ and $210$ (fixed basis) to $870$ and $612$, the largest single improvement. It does nothing about the non-skew term of Eq.~\eqref{eq:geq}, which is present in any frame, and with remaps switched off the high-$N_x$ run still fails (at 472 against 476). In all production studies of Chapter~\ref{ch:model} the width stays at the thermal width; the remap handles the drift.

The cheapest remedy is dissipation at the top. Lane B tested an order-two hypercollision of strength $\nu=1$ to 64 at $\vq\le0.03$ and $0.1$. Before the blow-up the observables do not depend on $\nu$ (spread $\le1.4\times10^{-6}$ at $t=600$), which is the evidence that the closure leaves the resolved physics untouched. But at $\vq=0.1$ every $\nu$ tested, $0.25$--$64$, fails before 1000; the closure postpones the loss of positivity and does not regularize. The fidelity cost against the ordinary benchmarks at $N=128$ is severe: the echo amplitude is $-81.7\,\%$ ($\nu=1$) and $-92.1\,\%$ ($\nu=2$) from the grid, and the echo time moves from $29.05$ to $28.55$ and $25.60$. At $N=512$ the cost falls to $-3.2$ and $-6.2\,\%$. A filter of the Hou--Li type (lane A: a width floor plus an exponential filter, SPECTRAX \#61) was never taken to a comparison with records; its branch and pull request remain open, and it is not claimed.

The field-scaled form of Eq.~\eqref{eq:closure} improves on the constant $\nu$ by scaling with the quantity that the growth bound of Eq.~\eqref{eq:growthbound} contains. It gives windows of $942$--$952$ at $\vq=0.03$ and $672$ at $0.1$ for $N_x=32$, with $c=0.5$, 1, 2 and $N_n=128$ agreeing under the $10\,\%$ rule. The next question was whether a strength $c(K,N_n)$ exists that restores stability at every refinement, and the closure study answered no: no $c$ in the ladder brings $\gamma\le10^{-3}$ at all 39 points with peak field $4\times10^{-3}$ and only at 23 of 39 with peak $10^{-3}$, and raising $c$ leaves the floor of Sec.~\ref{sec:frozen}. Gate Q1 failed; the decisive stage of scaling $c$ at $N_x=32$--128 and $N_n=64$--256 to $t=1000$ was therefore not run, and no option was added to SPECTRAX. A fixed strong closure on the real run (stage F, $N_x=64$, $N_n=128$, GPU) gives windows of $740$, $760$ and $720$ for $c=16$, 64 and 256, to be compared with 595 at $c=4$, and the fidelity to $c=4$ up to $t=500$ is $10^{-6}$ to $10^{-5}$ in $\Delta K_e$ and $10^{-7}$ in $\Wext$. A stronger closure therefore buys $120$ to $160$ time units at no cost in resolved physics and then stops; since $750$ is also where the grid reference ends, the saturation cannot be assigned to the AW floor alone.

The declared test of resolution as a remedy, $N_x\in\{64,128\}$ and $N_n\in\{128,256\}$ on a GPU at $\vq=0.1$ and $0.05$ with $c=4$, failed its gates G1 (agreement to 1000) and G4 (solver success to the end). All eight runs ended at the step budget of 200\,000 per 100-unit chunk, inside the instability, and the pairs agree to the end of the shorter run so the failure is of the resolution route, not of agreement. The certified windows are 328 and 440, G3 (ledger $\le10^{-6}$) held in every run ($\le1.7\times10^{-7}$), and G2 (positivity) failed in three.

The three candidate bases of the design note were prototyped in turn. Dai's stable Galerkin form is Galerkin with the same functions for trial and test. Written in AW coordinates in float64 it is unstable, with a frozen-field maximum real part of $10^{-2}$ to $1.1$ at $N_n=64$--128, Floquet exponents up to $1.0$, and an $L^2$ norm growing by $10^{88}$ by $t=1000$ at $N_n=64$. The cause is round-off, not the method: the trial space is the span of Hermite functions of width $a/\sqrt2$, so the method is the symmetric basis in disguise, and the change of basis from AW to orthonormal has condition number of order $2^N$ (matrix entries about $3\times10^{10}$ at $N_n=128$). In the orthonormal basis it passes the linear tests (largest real part $\le6\times10^{-15}$, OTSI within $0.96\,\%$ of Silin's relation), but it conserves only the moments of one parity of $N+1$, never number, momentum and energy together, and costs $1.4$ times the nonzeros and $2.5$ times the matrix-vector product. It was rejected. The plain symmetric basis (option B) has the neutral spectrum of Table~\ref{tab:aw}, but fails its nonlinear gates: for $\vq=0.1$, $\nu=0$, RK4 with $\mathrm dt=0.02$, every run ($N_x=16$ and 32, $N_n=64$--256) becomes nonfinite between $t=637.5$ and $737.5$, and $N_n=256$ dies earlier than 128 at $N_x=32$; the $N_x=16$ refinement pair agrees only to $t=634$ against a rejection threshold of 650; the energy ledger of the smallest run reaches $1.2\times10^{-3}\Wext$ before $t=600$; and the density and momentum, which are not exact in this basis, are off by $2.3\times10^{-3}$ and $1.1\times10^{-3}$ at $N_x=16$, $N_n=64$ (down to $5\times10^{-8}$ at $N_x=32$, $N_n=256$). The anti-symmetric square-root form of Issan et al.~\cite{Issan2024a,Issan2025} (in which $\sqrt f$ is expanded, so that $\sum C^2$ is the particle number) passed the linear stage: largest real part $\le4\times10^{-15}$ at 18 of 18 frozen and 12 of 12 Floquet points where the AW basis without closure is unstable at 18 of 18 and 11 of 12, Silin's rate to $0.25\,\%$, equality with the AW linearization to $2.5\times10^{-10}$, eigenvector condition 40 to 105 against $10^{19}$--$10^{29}$, and a cost of $1.5$ times the AW right-hand side. Number is conserved to $10^{-18}$, momentum and energy to the top-two-index fraction ($3\times10^{-3}$ at a fraction of $6\times10^{-4}$, and $10^{-12}$ to $10^{-14}$ for steep spectra). But the stage-two pilot, declared before the runs and with the gates of the SW study, departs from the grid at $t=391$ ($N_x=16$) and 498 ($N_x=32$), independently of $N_n$ (64 to 256) and of the step ($0.02$ against $0.005$), fails G3 and G4, and blows up within about ten time units after the top Hermite band fills; a diagnostic run at $N_x=64$ fails at 538. This is earlier than the SW basis (634 to 738). The question why the square-root $x$-discretization fails at a time that depends on $N_x$ but not on $N_n$, while the $f$-form tracks the grid (the two drivers agree to $10^{-9}$ until $t=250$), for instance through the $2/3$ mask acting on $g$ rather than on $f=g^2$, is open and was not run. That a basis proven bounded in $L^2$ fails earlier than the unbounded one reminds us that bounds on the semi-discrete linear operator do not imply nonlinear fidelity.

The closure Lawson factor of Sec.~\ref{sec:lawson} and the Hermite streaming factor are the two exponential integrators tried. The first gives $5.0\times$ fewer steps to $\wpe t=500$ and $4.3\times$ to 740 at $c=64$ at equal cost per step, with $\Delta K_e$ and $\Wext$ equal to the plain record to $5\times10^{-10}$ and the ledger at $8\times10^{-13}$; at $c=4$ the factor is $0.96$ and the declared gate (at least $3\times$ at both strengths) failed. It does not change the window. The streaming factor gave an identical step count at $c=64$, was $2.5\times$ slower in wall time, and at $c=4$ the run collapsed at 456.5 with 201\,927 steps in one chunk against 5444 for the closure-only control; the ordinary modulational run reached its budget of 400\,000 steps (932 s against 172 s). Section~\ref{sec:integrators} gives the reason these outcomes were expected.

\section{The certified window}\label{sec:windows}

If the instability cannot be removed it can be bounded. A \emph{certified window} is defined by checks declared before the runs: refinement pairs ($N_x$, $N_n$, closure strength) must agree to $10\,\%$ in $\Delta K_e$ and $\Wext$ (the running-maximum rule: the first $t>50$ at which the two differ by more than $10\,\%$ of the running maximum of the quantity), the electron distribution must stay positive (the kinetic energy of each species and the $k=0$ parallel temperature positive), the solver must reach the time, and the ledger of Eq.~\eqref{eq:ledger} must close. The window is the minimum of these times. It is a lower bound on where the Hermite results may be trusted, and a statement about the method and not about physics. An independent grid solver supplies a second reference but is itself limited: at $\vq=0.1$ it is certified to about 750 for $\Delta K_e$ and $\Wext$ and to about 640 for the spectrum, and it goes negative before the Hermite runs do, so it is not a better reference late in the run. Raising its velocity box from $32\vte$ to $48$ and $64$ moves edge contact only from $t=701.5$ to $754.5$, because the tail is accelerated past $64\vte$, and negativity at $645.5$ is set by $N_x$, not by the box.

The windows obtained are as follows. With no closure and $N_x=8$, 879 and 606.5 for $\vq=0.03$ and 0.1. With the field-scaled closure at $c=1$ and $N_x=16$ and 32, 942--952 and 672. With the strong closure $c=64$ integrated exactly by the Lawson factor and four members $(N_x,N_n)=(32,64),(64,64),(32,128),(64,128)$ plus $(32,64)$ at $c=32$, the windows used in the strong-drive study are 1000, 892 and 791.5 for $\vq=0.03$, 0.05 and 0.07: they are set by positivity loss of the $N_x=64$ members, that is by the truncation instability reaching the resolved range, not by disagreement of the refinement pairs. At $\vq=0.1$ there is no window beyond about 760 (760 at $c=64$ by positivity, 579.5 at $c=4$ by the step budget). No drive reaches $t=1000$ at $\vq=0.03$--$0.1$ with $N_x$-converged results other than the three values just given for $c=64$, and the collisionless windows quoted for the paper at the HHS value 0.03 are therefore 1000. Within the windows the work follows the exact linear resonant law to $0.68\,\%$ at $\vq=0.07$ (Chapter~\ref{ch:model}, Eq.~\eqref{eq:linlaw}), the four members agree to $2\times10^{-4}$, and the ledger closes to $\le7\times10^{-12}$.

\section{The physical oscillating two-stream instability}\label{sec:otsi}

The strong drive also excites a real instability, and it is important not to confuse the two. A uniform dipole pump of frequency $\omega_0$ gives the electrons an oscillating excursion $r(t)$. In the frame of the oscillating electrons the ions stream through, and finite-$k$ perturbations couple through Bessel factors, $e^{\ii kr\sin\omega_0t}=\sum_lJ_l(kr)e^{\ii l\omega_0t}$, to sidebands $\omega+l\omega_0$. Silin~\cite{Silin1965} writes the dispersion relation as the determinant of a sideband matrix, which for Maxwellian susceptibilities $\chi_s$ we evaluate as $\det[\mathrm{diag}(1+\chi_i)+B\,\mathrm{diag}(\chi_e)B^T]=0$ with $B_{lj}=J_{l-j}(kr)$ and $|l|\le4$, taking the ions unpumped. Nishikawa's coupled-mode theory~\cite{Nishikawa1968a,Nishikawa1968b} gives the same classification and a cube-root dependence on pump power, but for a fluid plasma with phenomenological damping, so it serves as classification and scaling and not as a quantitative reference. The linear pump-frame Hermite model with a constant pump $E_0\cos\omega_0t$ has Floquet exponent from the one-period monodromy matrix; it agrees with the Silin determinant to $1.3\,\%$ at all 22 unstable points ($v_{\rm os}/\vte=0.3$--10, $k=2$--$30\,k_1$), and the modes are purely growing in the ion frame, which identifies the OTSI and excludes the decay or parametric-decay instability. With fixed ions every rate is below $2\times10^{-12}$, so there is no electron-only mechanism. The dependence on the pump frequency, at $k=12k_1$, $v_{\rm os}=3\vte$, is rate $0.0067$, $0.010$, $0.012$, $0.0097$, $0.0077$, $0$, $0$ for $\omega_0/\wpe=0.9$, 0.97, 1, 1.003, 1.01, 1.05, 1.2, so that the instability exists only on the OTSI side of resonance (the decay side is stable at $T_i=T_e$); the rate is insensitive to $T_i/T_e$ and scales as $m_i^{-0.29}$, close to the weak-field $(m_e/m_i)^{1/3}$ of Silin's analysis. Integrating the Floquet rate along the secular excursion of the resonant drive, $2\int\gamma(k,v_{\rm rel}(t))\,\dd t$, predicts base-ten logarithms of the gain of $|E_k|^2$ at $\wpe t=700$ of $2.61$, $6.41$, $9.90$ for $k=2,5,8$ times $k_1$, against $1.80$, $5.78$, $9.31$ in the linear Hermite run, uniformly $0.6$--$0.8$ decades high, so that the high-$k$ growth of the resonant run is the quasi-static OTSI. The relation is established physics in the sense of Silin and Nishikawa; what is new is the regime, a resonant pump with a secularly growing excursion and $T_i=T_e$.

The linear Hermite model does not invent the growth. Hermite orders 32, 64 and 128 agree with each other for every mode, they follow the exact Volterra solution until the first $10\,\%$ difference at $\wpe t=405$ ($k_1$), 323, 292, 273, 270, 276, 317 and 364 ($k_8$), and end below it by factors of 1.3 to 4. Hence the physical amplification is understated by the Hermite model, and the numerical instability of this chapter, whose growth rate increases with $N_n$ (Table~\ref{tab:aw}), is a different phenomenon superimposed. Two diagnostics distinguish them in practice. The numerical mode sits in the top third of the Hermite ladder at the largest retained $k$, the physical OTSI sits at low Hermite index and rising $k$; and the numerical mode ends the run sooner when the resolution is increased, which a resolved physical mode does not. In the nonrelativistic strong-drive study the onset of the cascade (fluctuating-field amplitude reaching $2\times10^3E_0^2$) is at $\wpe t=783.5$ for $N_x=32$ and $862$ for the higher-resolution members, while the run ends by positivity loss afterwards; within the windows the instability has not extracted a resolvable fraction of the work. The statement that no saturation appears before the cascade begins is the main limitation of the bound derived in the paper and says nothing about the evolution after it.

\section{What is open}\label{sec:open-limits}

Three questions remain. A time-domain, nonlinear stable discretization that conserves mass, momentum and energy and preserves the Gaussian weight is not available: SW and its square-root form give up exact moments and fail earlier, AW gives up stability. The failure of the square-root form at a time dependent on $N_x$ but not $N_n$ has not been diagnosed. And the frozen-field spectrum has been analysed for a single species in one velocity dimension and for the electric force, while the magnetic variant (Weibel, Chapter~\ref{ch:verification}) is driven by $\vv\times\BB$ and requires its own closure. The practical consequence is the certified-window discipline described above: results at strong drive are reported with their windows, the figures show only certified intervals, and claims past a window are not made.
